\section{Example: Simple Build-In Optimization} \label{ex5}

{~}\\[-0.5cm]
A very common situation is, that you want to tune input parameters
with respect to a figure of merit. \W has build in two simple
procedures to accomplish that. The first is an iterativ estimation
for the case of one parameter to adjust, the second a Monte-Carlo for more
than one parameter.

Consider example \ref{ex4} For the given gap, we want to find the
shift such, that the vertical and horizontal field of the undulator are
the same, i.e. the device acts as an circular undulator.

In a terminal, we enter:\\

\cboxs{
\senv\\ \grau
python3 ../python/wave\_example\_5.py\\
cp ../input/wave\_example\_5.in wave.in\\
wave\\
}

We see in \woe, that the max. vertical and horizontal fields \scap{BYMAX} and
\scap{BZMAX} differ. Now, we
adjust the shift parameter of the undulator. Enter:\\

\cboxs{cp ../input/wave\_example\_5.adjust wave.adjust\\
cp ../python/wave\_adjust\_example\_5.py wave\_adjust.py\\
}

Have a look on \gtis{wave.adjust} (Fig.\ref{ex5fig2}).
When a file like this is passed as an argument to \We, \W will adjust the
parameter \scap{URSHIFT(1)} of \wi by a Monte-Carlo search. Of course,
this also works with other filenames and variables of \wie.  During
the optimization
procedure \W executes the script \gtis{wave\_adjust.sh}, which calls
another instance of \W and \gtise{wave\_adjust.py}.\\

\begin{figure}[htb]
    \centering
    \includegraphics[width=0.8\textwidth]
    {/home/scheer/wav/doc/pp/wave_example_5_wave_adjust.png}
    \caption{Input file for the optimization procedure}
    \label{ex5fig1}
\end{figure}

After the run of this instance,
the command given in the last line of \gtis{wave.adjust} is executed by
the master process.
In our example, this is the python script  \gtis{wave\_adjust.py}.
What does this script (Fig.\ref{ex5fig2}) do? \\

\begin{figure}[htb]
    \centering
    \includegraphics[width=0.8\textwidth]
    {/home/scheer/wav/doc/pp/wave_example_5_wave_adjust_py.png}
    \caption{Example of a program to calculate the figure of
    merit}
    \label{ex5fig2}
\end{figure}

It opens \woe, and looks for the lines containing the keywords like
\scap{BYmin}. When a matching line is found,
the program reads from the line the values of  \scap{By} or \scap{Bz},
respectively,
and calculates the \gtis{$\chi^2$} to be minimized. The resulting \gtis{$\chi^2$}
is written to \gtis{wave\_adjust.dat} and evaluated by the
master process.

Let's try. Enter:

\cboxs{cp wave.in wave\_adjust.in\\
wave wave.adjust \\
}


\W loops and writes \gtis{wave\_adjust.out}.
The data in this file show a minimum for the shift of about \gtis{19mm}
(Fig.\ref{ex5fig3}). This minimum value found is also used to re-run
WAVE.

%\newpage
Start the plotting GUI by entering:

\cboxs{wplot}

and enter in the python shell:\\

\cboxs{adj = ncread("adj","sh:chi2","wave\_adjust.out")\\
nplot(adj,"sh:chi2")\\
txyz("wave\_adjust.out","Shift[mm]","Chi2")}

\begin{figure}[htb]
    \centering
    %\framebox(12,9){
    \includegraphics[width=0.7\textwidth]
    {/home/scheer/wav/doc/pp/wave_adjust_mc.pdf}
    %}
    \caption{Random variation of URSHIFT(1)}
    \label{ex5fig3}
\end{figure}

The plot you get, might differ from that shown here, due to different
seeds of the Monte-Carlo search.

To get a precise value,
we edit \gtis{wave.adjust}, and set the number of parameters to \gtis{-1}
and the boundaries to \gtis{18.0} and \gtis{20.0}.
Then we repeat

\cboxs{cp wave.in wave\_adjust.in\\
wave wave.adjust\\
}

and find now in \wi the best value of about \\

\scap{URSHIFT(1)= 18.6}\\

Since we have required a precision for the minimum
\gtis{$\chi^2 \le 10^{-6}$}, the ratio of the vertical and horizontal field
now agrees on the scale of \gtis{$\sqrt{10^{-6}}$}.

The effort in this example might appear quite high for such a simple
probleme, which could be solve sufficently by a quick interpolation.
However, it shows the usage of scripts and the advantage of command
line driven optimization procedures.
